% !TeX root = ../初等数论笔记.tex
% 由主文档通过 \input 载入；请编译主文档。
\chapter{绪言与课程路线}
\label{ch:course-outline}

\emph{初等数论（elementary number theory）}研究整数的整除、分解、同余及整数方程。本书面向本科生，以中学代数和基本逻辑为起点，不预设抽象代数。内容围绕三个问题展开：整数如何分解，取余运算保留哪些结构，以及整数方程的解如何判定与构造。

全书分为四个阶段。第\ref{ch:introduction}至\ref{ch:primes}章建立证明方法、整除理论与唯一分解；第\ref{ch:congruences}至\ref{ch:quadratic-residues}章依次讨论同余、中国剩余定理、算术函数、幂同余、原根与二次互反律；第\ref{ch:diophantine}至\ref{ch:pell}章讨论整数方程、连分数和 Pell 方程；第\ref{ch:prime-distribution}、\ref{ch:algorithms}章选讲素数分布与算法应用。前一阶段建立的工具会在后一阶段反复使用：贝祖等式产生模逆元，模逆元支撑中国剩余定理，素因子分解支撑算术函数，乘法阶和平方剩余又把同余方法带入整数方程。

正文突出定义、定理、证明与例题。核心推导尽量自足；较长的拓展定理在引用处明确标注不证，不以证明思路代替完整证明。章末习题用于检验概念、训练证明，附录提供代表题的答案与综合解答。计算结果用于核对例子，不能代替一般命题的证明。

学习时可以分两遍阅读。第一遍掌握每个定义的对象与条件，亲手完成例题，理解主要定理解决什么问题；第二遍追踪证明中每一步的依据，尤其注意互素、素数、非平方等条件在哪里使用。原根分类、二次互反律的格点证明、平方根连分数的周期证明较长，初读时可先把握结论和应用，再返回细读。带有“选学”的内容不作为基础习题的默认前提；素数分布中涉及级数和极限的部分需要微积分基础，其余主线以整数运算和初等代数为主。

数论题目的答案通常有几个层次。“找到一组解”只是存在性证据；“写出一族解”还要检查参数条件；“求出全部解”则必须证明任何解都能落入这族参数。阅读例题时，应分别辨认构造与完备性论证。各章习题前四题用于巩固本章主线，后续题按概念辨析、方法迁移和综合证明逐步展开；附录给出代表题的答案与后四题的关键步骤，宜在独立尝试后使用。

课程范围可与 MIT 本科课程资料\cite{mit-syllabus,mit-notes}对照，计算专题可进一步阅读 Stein 的教材\cite{stein-book}。参考资料中的代数和椭圆曲线部分不作为本书先修要求。

